\documentclass[10pt,a4paper]{amsart}
\usepackage[margin=19mm]{geometry}
\usepackage{amsmath,amssymb,mathtools}
\usepackage[scr=rsfs,cal=euler]{mathalfa}
\usepackage{xfrac}
\usepackage[unicode,hidelinks,bookmarks=false]{hyperref}
\newcommand{\ie}{i.e.\ }
\newcommand{\cf}{cf.\ }
\newcommand{\resp}{respectively\ }
\newcommand{\Subsection}{\subsection}
\providecommand{\nobreakdash}{\nobreak\hbox{-}}
\setlength{\emergencystretch}{2em}
\multlinegap0pt
\usepackage{bm}
\usepackage{bbm}
\usepackage{dsfont}
\usepackage{xspace}
\usepackage{cancel}
\usepackage{mathtools}
\usepackage{amsthm}
\makeatletter
\@ifundefined{lemma}{\newtheorem{lemma}{Lemma}}{}
\makeatother
\newcommand{\eqdef}{\ensuremath{\stackrel{\mbox{\upshape\tiny def.}}{=}}}
\title[Bridging the Gap Between Approximation and Learning via Opti]{Bridging the Gap Between Approximation and Learning via Optimal Approximation by ReLU MLPs of Maximal Regularity}
\author{Ruiyang Hong, Anastasis Kratsios}
\date{Source: arXiv:2409.12335v4 (CC BY 4.0)}
\begin{document}
\maketitle

\noindent\emph{Source and licence.} Bridging the Gap Between Approximation and Learning via Optimal Approximation by ReLU MLPs of Maximal Regularity. Version arXiv:2409.12335v4 (CC BY 4.0), distributed under the Creative Commons Attribution 4.0 International licence, as recorded in the arXiv licence metadata for this exact version and shown on the abstract page https://arxiv.org/abs/2409.12335v4. Full rights notice: licenses/arxiv-2409.12335-CC-BY-4.0.txt.\par\smallskip

\noindent\textit{Editorial scope.} Counted unit: the Lemma whose source label is \texttt{lem:continuous\_piecewise\_linear\_approximators\_preserve\_regularity} in the article (source0.tex lines 1645--1652, source label \texttt{lem:continuous\_piecewise\_linear\_approximators\_preserve\_regularity}), delivered together with the article's own complete original proof of it (source0.tex lines 1654--1795, one \texttt{proof} environment of 142 lines). No mathematics is written, repaired or re-derived: statement and proof are the article's own text line for line. Required context retained uncounted: none. Every definition, display and label of those lines is retained unchanged. Omitted from this package and not redistributed: 1--1644, 1653--1653, 1796--4778. Nothing inside a retained excerpt is omitted or altered: every retained body is delivered exactly as the article's lines are written, so no omission receipt of any kind accompanies this package. Citations: the retained excerpts carry no citation command at all, so no bibliography entry of the article is transcribed and no reference is redistributed. Reference adaptation: none was needed and none was made; every reference inside the delivered unit resolves inside it, so no unresolved reference or citation is produced. Printed slips kept literal: no printed word, symbol, hypothesis or number of the counted statement or of its proof was altered. Layout and class: the article's own class is \texttt{article} with packages including geometry, inputenc, float, graphicx, amsthm, amsmath, amssymb, mathabx, mathrsfs, xcolor, array, multirow, caption, amsfonts; the wrapper is the lane's pinned \texttt{amsart} editorial wrapper at 10pt on A4, which restates the theorem-like environments the retained text needs and adds the article's own macro definitions that the retained text needs (\texttt{\textbackslash eqdef}), with extra wrapper packages (bm, bbm, dsfont, xspace, cancel, mathtools). The delivered numbering is the unit's own. Assets: no retained span contains a figure environment, a graphics-inclusion command, a PDF-page inclusion, a table or a TikZ picture; no image file is copied into this package, the package declares no asset dependency and no image file is redistributed with it. Licence: version arXiv:2409.12335v4 of this article is distributed under the Creative Commons Attribution 4.0 International licence, as recorded in the arXiv licence metadata for this exact version and shown on the abstract page https://arxiv.org/abs/2409.12335v4. A full rights notice is delivered at licenses/arxiv-2409.12335-CC-BY-4.0.txt. Characters: every retained excerpt is pure ASCII, so the delivered engine, XeLaTeX with its default Latin Modern text font, reports no missing character.
\begin{lemma}[Continuous piecewise linear approximators preserve regularity, and $L^{\infty}$ error estimate]
\label{lem:continuous_piecewise_linear_approximators_preserve_regularity}
    Let $f$ be a function from $[a,b]$ to $\mathbb{R}$ with modulus of regularity $\omega$, and let $\Phi$ be a continuous piecewise linear function on $[a,b]$ that passes through the points $\{(x_i,f(x_i))\}_{i=0}^n$ for $a=x_0<x_1<x_2<\cdots<x_n=b$, and is linear on $[x_0,x_1],[x_1,x_2],\cdots,[x_{n-1},x_n]$. Then, we have the error estimate 
    $$\| f-\Phi \|_{L^{\infty}([a,b])} \le \omega \left(
    \max\limits_{i=0,1,\cdots,n-1}{\frac{|x_{i+1}-x_i|}{2}}
    \right)$$
    Moreover, $\omega$ is a modulus of regularity of $\Phi$. 
\end{lemma}

\begin{proof}

We first show that $\Phi$ has at least the same regularity as $f$. Take any $x,y\in[a,b]$, and assume without loss of generality that $x\le y$. If $x,y$ lie in the same interval in $[x_0,x_1],[x_1,x_2],\cdots,[x_{n-1},x_n]$, then $x,y\in[x_m,x_{m+1}]$ for some index $m$, by concavity of $\omega$, 
\begin{align*}
    \omega(y-x) 
    &= 
    \omega\left(
        \left( 1 - \frac{y-x}{x_{m+1}-x_m} \right) \cdot 0
        +
        \frac{y-x}{x_{m+1}-x_m} \cdot (x_{m+1}-x_m)
    \right)
    \\ 
    &\ge 
    \left( 1 - \frac{y-x}{x_{m+1}-x_m} \right) \omega(0)
    +
    \frac{y-x}{x_{m+1}-x_m} \omega(x_{m+1}-x_m)
    \\
    &\ge 
    \frac{y-x}{x_{m+1}-x_m} \omega(x_{m+1}-x_m)
    \\
    &\ge 
    \frac{y-x}{x_{m+1}-x_m} | f(x_{m+1}) - f(x_m) |
    \\
    &= 
    \frac{y-x}{x_{m+1}-x_m} | \Phi(x_{m+1}) - \Phi(x_m) |
    \\
    &=
    | \Phi(y) - \Phi(x) |
\end{align*}
where the last equality follows from the fact that $\Phi$ is linear on $[x_m,x_{m+1}]$. Otherwise, $x,y$ lie in different intervals in $[x_0,x_1],[x_1,x_2],\cdots,[x_{n-1},x_n]$, then $x_p \le x \le x_{p+1} \le x_q \le y \le x_{q+1}$ for some indices $p,q$, and let 
$$A \eqdef \frac{x-x_p}{x_{p+1}-x_p}, \quad
B \eqdef \frac{x_{q+1}-y}{x_{q+1}-x_q}$$
so $0 \le A,B \le 1$. Without loss of generality, assume that $A\le B$. Then, 
$$\Phi(x) = (1-A) \Phi(x_p) + A \Phi(x_{p+1}), \quad
  \Phi(y) = B \Phi(x_q) + (1-B) \Phi(x_{q+1})$$
Since 
\begin{align*}
    y-x
    &=
    \left(\frac{x_{q+1}-y}{x_{q+1}-x_q} - \frac{x-x_p}{x_{p+1}-x_p}\right)
    (x_q-x_p)
    + \left(1-\frac{x_{q+1}-y}{x_{q+1}-x_q}\right)
    (x_{q+1}-x_p)  \\
    &+ \frac{x-x_p}{x_{p+1}-x_p}(x_q - x_{p+1})  \\
    &= (B-A)(x_q-x_p) + (1-B)(x_{q+1}-x_p) + A(x_q - x_{p+1})
\end{align*}
with $B-A\ge 0, 1-B\ge 0, A\ge 0$ and $(B-A)+(1-B)+A=1$, by concavity of $\omega$, we have 
\begin{align*}
    \omega(y-x) 
    &= \omega( (B-A)(x_q-x_p) + (1-B)(x_{q+1}-x_p) + A(x_q - x_{p+1}) )  \\
    &\ge (B-A) \omega(x_q-x_p)
    + (1-B) \omega (x_{q+1}-x_p)
    + A \omega (x_q - x_{p+1})
    \\
    &\ge (B-A) |\Phi(x_q) - \Phi(x_p)|
    + (1-B) |\Phi(x_{q+1}) - \Phi(x_p)|
    + A |\Phi(x_q) - \Phi(x_{p+1})|
    \\
    &\ge (B-A) (\Phi(x_q) - \Phi(x_p))
    + (1-B) (\Phi(x_{q+1}) - \Phi(x_p))
    + A (\Phi(x_q) - \Phi(x_{p+1}))
    \\
    &= (B \Phi(x_q) + (1-B) \Phi(x_{q+1})) 
    - ((1-A) \Phi(x_p) + A \Phi(x_{p+1}))
    \\
    &= \Phi(y) - \Phi(x)
\end{align*}
Similarly, we can show that $\omega(y-x) \ge \Phi(x) - \Phi(y)$, thus 
$$\omega(|y-x|) = \omega(y-x) \ge |\Phi(y) - \Phi(x)|$$
Therefore, we always have $|\Phi(y) - \Phi(x)| \le \omega(|y-x|)$ for any $x,y\in[a,b]$, thus $\omega$ is a modulus of regularity of $\Phi$.

Now we prove the upper bound for the $L^{\infty}$ error. Take any $x\in[a,b]$, then $x\in[x_k,x_{k+1}]$ for some index $k$. For convenience, let $L \eqdef x_{k+1}-x_k$, then we have 
\allowdisplaybreaks
\begin{align*}
    |f(x) - \Phi(x)|
    &= 
    \left| f(x) - 
        \left(
            \frac{x_{k+1}-x}{L} \Phi(x_k)
            + \frac{x-x_k}{L} \Phi(x_{k+1})
        \right) 
    \right|
    \\
    &=\left| f(x) - \left(
        \frac{x_{k+1}-x}{L} f(x_k)
        + \frac{x-x_k}{L} f(x_{k+1})
    \right) \right|
    \\
    &=
    \left|
    \frac{x_{k+1}-x}{L} (f(x) - f(x_k))
    +
    \frac{x-x_k}{L} (f(x) - f(x_{k+1}))
    \right|
    \\
    &\le 
    \frac{x_{k+1}-x}{L} |f(x) - f(x_k)|
    +
    \frac{x-x_k}{L} |f(x) - f(x_{k+1})|
    \\
    &\le
    \frac{x_{k+1}-x}{L} \omega(x - x_k)
    +
    \frac{x-x_k}{L} \omega(x_{k+1} - x)
    \\
    &\le
    \omega \left(
        \frac{x_{k+1}-x}{L} (x - x_k)
        +
        \frac{x-x_k}{L} (x_{k+1} - x)
    \right)
    \\
    &= 
    \omega \left(
        \frac{(L-(x - x_k))(x - x_k)}{L} 
        +
        \frac{(x - x_k)(L-(x - x_k))}{L} 
    \right)
    \\
    &\le
    \omega \left(
        \frac{\left(\frac{L^2}{4}\right)}{L} 
        +
        \frac{\left(\frac{L^2}{4}\right)}{L} 
    \right)
    \\
    &=
    \omega\left( \frac{L}{2} \right)
    \\
    &\le
    \omega \left(
    \max\limits_{i=0,1,\cdots,n-1}{\frac{|x_{i+1}-x_i|}{2}}
    \right)
\end{align*}
Since $x\in[a,b]$ was chosen arbitrarily, we conclude that 
$$\| f-\Phi \|_{L^{\infty}([a,b])} \le \omega \left(
    \max\limits_{i=0,1,\cdots,n-1}{\frac{|x_{i+1}-x_i|}{2}}
\right).
% \qedhere
$$
This completes our proof.
\end{proof}

\end{document}
